\documentclass[11pt,a4paper]{article}

% ---------------- PACKAGES ----------------
\usepackage[utf8]{inputenc}
\usepackage[T1]{fontenc}
\usepackage[french]{babel}
\usepackage{amsmath, amssymb}
\usepackage{siunitx}
\usepackage{geometry}
\usepackage{multicol}
\usepackage{tcolorbox}
\usepackage{tikz}
\usepackage{xcolor}
\usepackage{fancyhdr}

\geometry{margin=2cm}
\pagestyle{fancy}
\fancyhf{}
\lhead{\textbf{Puissance et Énergie}}
\rhead{\thepage}

% ---------------- COULEURS ----------------
\definecolor{energyblue}{RGB}{30,80,160}
\definecolor{powerred}{RGB}{180,50,50}
\definecolor{warningred}{RGB}{200,30,30}
\definecolor{graybg}{RGB}{245,245,245}

% ---------------- TCOLORBOX ----------------
\tcbset{
  sharp corners,
  colback=white,
  colframe=black!60,
  boxrule=0.8pt,
  left=6pt, right=6pt, top=4pt, bottom=4pt
}

\newcommand{\theorybox}[2]{
\begin{tcolorbox}[title=\textbf{#1}]
#2
\end{tcolorbox}
}

% ---------------- GRILLE EXERCICES ----------------
\newcommand{\fullgrid}{
\vfill
\begin{center}
\begin{tikzpicture}[scale=0.9]
\draw[step=0.5cm,gray!30,very thin] (0,0) grid (18,22);
\end{tikzpicture}
\end{center}
\vfill
}

\newcommand{\threequartergrid}{
\vfill
\begin{center}
\begin{tikzpicture}[scale=0.9]
\draw[step=0.5cm,gray!30,very thin] (0,0) grid (18,18);
\end{tikzpicture}
\end{center}
\vfill
}

\begin{document}

% =================================================
% PAGE DE TITRE
% =================================================

\begin{center}
{\LARGE \textbf{Puissance et Énergie}}\\[0.3cm]
{\large Rappels de base, unités SI, conversions et exercices}\\[0.5cm]
\end{center}

% =================================================
% PARTIE I : RAPPELS THÉORIQUES
% =================================================

\section*{I - Rappels fondamentaux}

% -------------------------------------------------
\theorybox{A - Les 7 unités de base du Système International (SI)}{
Le Système International repose sur \textbf{7 unités fondamentales}.  
Toutes les autres unités sont construites à partir d’elles.

\begin{center}
\begin{tabular}{|c|c|c|}
\hline
\textbf{Grandeur} & \textbf{Unité} & \textbf{Symbole} \\
\hline
Longueur & mètre & m \\
Masse & kilogramme & kg \\
Temps & seconde & s \\
Courant électrique & ampère & A \\
Température & kelvin & K \\
Quantité de matière & mole & mol \\
Intensité lumineuse & candela & cd \\
\hline
\end{tabular}
\end{center}
}

% -------------------------------------------------
\theorybox{B - Unités dérivées courantes}{
Les unités dérivées sont des combinaisons des unités SI de base.

\begin{itemize}
\item Distance : cm, km (mais \textbf{le mètre est l’unité SI})
\item Masse : g, tonne (mais \textbf{le kilogramme est l’unité SI})
\item Volume : litre (L), m$^3$ (\textbf{toujours ramener en m$^3$})
\item Vitesse : m/s, km/h
\item Énergie : joule (J)
\item Température : °C (attention : le SI est le kelvin)
\item Charge électrique : coulomb (C) / ampère-heure (Ah)
\item Tension électrique : volt (V)

\end{itemize}

Exemples de conversion et relations :
\[
1\,\si{L} = 10^{-3}\,\si{m^3}
\qquad
1\,\si{km/h} = \frac{1000}{3600}\,\si{m/s}
\]

Relations électriques :
\[
1\,\si{C} = 1\,\si{A} \cdot 1\,\si{s}
\qquad
1\,\si{V} = 1\,\frac{\si{J}}{\si{C}} = 1\,\frac{kg \cdot m^2}{s^3 \cdot A}
\qquad
1\,\si{Ah} = 3600\,\si{C}
\]
}

% -------------------------------------------------
\theorybox{C - RÈGLE ABSOLUE - UNITÉS SI}{
\centering
\textbf{\Large TOUT DOIT ÊTRE EN UNITÉS SI AVANT D’UTILISER UNE FORMULE}

\vspace{0.2cm}

Pas de grammes.  
Pas de tonnes.  
Pas de kilomètres.  
Pas de centimètres.  

\textbf{kg, m, s, A, K.}

Sinon, \textbf{le résultat est faux}, même si le calcul est juste.
}

% -------------------------------------------------
\theorybox{D - Conversion par déplacement de la virgule (longueurs, masses)}{
Pour les longueurs et les masses, on utilise une \textbf{table de conversion décimale}.

Exemple : Conversion de $1470\,\si{mm}$ en mètres
\textbf{Table de conversion :}

\begin{center}
\begin{tabular}{c|c|c|c|c|c|c}
km & hm & dam & m & dm & cm & mm \\
\hline
   &    &     & 1 & 4  & 7  & 0  \\
\end{tabular}
\end{center}

Chaque déplacement d’une colonne vers la gauche correspond à une division par 10.

\vspace{0.2cm}

\textbf{Écriture décimale avec déplacement de la virgule :}

\[
1470\,\si{mm}
\;\longrightarrow\;
147.\!0\,\si{cm}
\;\longrightarrow\;
14.\!70\,\si{dm}
\;\longrightarrow\;
1.\!470\,\si{m}
\]

\vspace{0.2cm}

\textbf{Résultat final :}

\[
1470\,\si{mm} = 1.47\,\si{m}
\]




Chaque déplacement vers la droite multiplie par 10.  
Chaque déplacement vers la gauche divise par 10.

Exemple :
\[
250\,\si{cm} = 2.5\,\si{m}
\qquad
3.2\,\si{km} = 3200\,\si{m}
\]
}

% -------------------------------------------------
\theorybox{E - Conversion des surfaces et volumes (tableau dimensionnel)}{

\textbf{Principe visuel :}

\begin{itemize}
\item Surface $\rightarrow$ deux dimensions $\rightarrow$ \textbf{deux zéros}
\item Volume $\rightarrow$ trois dimensions $\rightarrow$ \textbf{trois zéros}
\end{itemize}

\vspace{0.2cm}

% ================= SURFACES =================
\textbf{Exemple - Conversion de $25\,\si{cm^2}$ en $\si{m^2}$}

\begin{center}
\textbf{Tableau des surfaces}
\begin{tabular}{c|c|c|c|c|c|c}
km$^2$ & hm$^2$ & dam$^2$ & m$^2$ & dm$^2$ & cm$^2$ & mm$^2$ \\
\hline
       &        &         & 0 0 & 2 5 &       &        \\
\end{tabular}
\end{center}

Lecture du tableau :
\[
25\,\si{cm^2} = 0.0025\,\si{m^2}
\]

\[
25\,\si{cm^2} = 2.5 \times 10^{-3}\,\si{m^2}
\]

\vspace{0.4cm}

% ================= VOLUMES =================
\textbf{Exemple - Conversion de $4\,\si{cm^3}$ en $\si{m^3}$}

\begin{center}
\textbf{Tableau des volumes}
\begin{tabular}{c|c|c|c|c|c|c}
km$^3$ & hm$^3$ & dam$^3$ & m$^3$ & dm$^3$ & cm$^3$ & mm$^3$ \\
\hline
       &        &         & 0 0 0 & 0 0 4 &        &        \\
\end{tabular}
\end{center}

Lecture du tableau :
\[
4\,\si{cm^3} = 0.000004\,\si{m^3}
\]

\[
4\,\si{cm^3} = 4 \times 10^{-6}\,\si{m^3}
\]

\vspace{0.2cm}

\textbf{Message clé :}

Les zéros ne sont pas décoratifs.  
Ils codent le \textbf{nombre de dimensions}.

}



% -------------------------------------------------
\theorybox{F - Conversion par proportion (cas général)}{
Quand la conversion n’est pas décimale, on utilise une \textbf{proportion}.

Exemples :

\[
1\,\text{lb} = 453\,\si{g}
\Rightarrow
5\,\text{lb} = 5 \times 0.453 = 2.27\,\si{kg}
\]

\[
1\,\si{km/h} = \frac{1000}{3600}\,\si{m/s}
\Rightarrow
120\,\si{km/h} = 33.3\,\si{m/s}
\]
}



% -------------------------------------------------
\theorybox{G - Définition de l’énergie (intuition physique)}{

L’énergie ne se voit pas directement, mais on la \textbf{ressent} lorsqu’un effort est nécessaire.

\vspace{0.2cm}

\textbf{Example : Le vélo et la colline}

\vspace{0.2cm}

Un cycliste roule sur du plat, puis arrive au pied d’une colline.

Pour monter, il doit \textbf{fournir un effort} :  
ses jambes exercent une force pour élever le vélo et le cycliste contre la gravité.

Cette énergie fournie ne disparaît pas.  
Elle est stockée sous forme d’\textbf{énergie potentielle de pesanteur} :

\[
E_p = m g h
\]

où :
\begin{itemize}
\item $m$ est la masse du système (cycliste + vélo) en [Kg],
\item $g$ l’accélération de la pesanteur en [m/s²],
\item $h$ la hauteur gagnée en [m].
\end{itemize}

Plus la colline est haute ou plus la masse est grande, plus l’énergie stockée est grande.  
C’est exactement ce que ressent le cycliste : \textbf{monter coûte de l’énergie}.

\vspace{0.3cm}

Arrivé en haut, le cycliste s’arrête un instant.  
Le vélo ne bouge pas, mais l’énergie est toujours là.

\vspace{0.2cm}

\textbf{Descente}

\vspace{0.2cm}

Lorsque le cycliste se laisse redescendre sans pédaler, le vélo \textbf{accélère}.
Il n’y a pourtant aucun moteur.

L’énergie potentielle stockée se transforme alors en \textbf{énergie cinétique},
c’est-à-dire en énergie de mouvement :

\[
E_c = \frac{1}{2} m v^2
\]

Plus la vitesse augmente, plus l’énergie cinétique est grande.

\vspace{0.3cm}

\textbf{Idée clé}

\vspace{0.1cm}

L’énergie ne disparaît pas :  
elle \textbf{change de forme}.

\[
E_p \longrightarrow E_c
\]

Sans frottements, toute l’énergie potentielle gagnée à la montée
serait retrouvée sous forme d’énergie cinétique à la descente.

\vspace{0.3cm}

\textbf{Unité}

\vspace{0.1cm}

Toutes ces formes d’énergie se mesurent en \textbf{joules} :

\[
[J] = kg \cdot m^2 \cdot s^{-2}
\]

C’est la même unité pour l’effort du cycliste, monter la colline
et prendre de la vitesse à la descente.

}


% -------------------------------------------------
\theorybox{H - Définition de la puissance (intuition physique)}{

La puissance mesure la \textbf{vitesse à laquelle l’énergie est transférée}.
Ce n’est pas une quantité d’énergie, mais un \textbf{débit d’énergie}.

\vspace{0.2cm}

\[
P = \frac{E}{t}
\qquad\Longleftrightarrow\qquad
1\,\si{W} = 1\,\si{J/s}
\]

\vspace{0.3cm}

\textbf{Exemple : l’ascenseur}

\vspace{0.2cm}

Un ascenseur de masse $m$ doit monter à une hauteur $h$.

Pour cela, il faut lui fournir une énergie de pesanteur :

\[
E = m g h
\]

Cette énergie est \textbf{imposée par la physique} :  
elle ne dépend ni du moteur, ni du temps de montée.

\vspace{0.3cm}

\textbf{Deux moteurs, même ascenseur}

\vspace{0.2cm}

\begin{itemize}
\item Un moteur peu puissant met beaucoup de temps à monter.
\item Un moteur très puissant monte rapidement.
\end{itemize}

Dans les deux cas, l’énergie totale fournie est la même.  
Ce qui change, c’est \textbf{la durée du transfert}.

\vspace{0.3cm}

\textbf{Puissance comme débit}

\vspace{0.2cm}

La puissance indique \textbf{combien de joules sont fournis chaque seconde}.

\begin{itemize}
\item $500\,\si{W}$ signifie $500\,\si{J}$ par seconde.
\item $2000\,\si{W}$ signifie $2000\,\si{J}$ par seconde.
\end{itemize}

Un moteur plus puissant fournit la même énergie,
mais avec un \textbf{débit plus élevé}.

\vspace{0.3cm}

\textbf{Conséquence directe}

\vspace{0.1cm}

Si la puissance est multipliée par 4,  
le temps de montée est divisé par 4 :

\[
t = \frac{E}{P}
\]

\vspace{0.3cm}

\textbf{Lecture en unités SI}

\vspace{0.1cm}

L’unité de la puissance est le watt :

\[
[W] = \frac{J}{s}
\]

Or, l’unité du joule est :

\[
[J] = kg \cdot m^2 \cdot s^{-2}
\]

Donc :

\[
[W] = kg \cdot m^2 \cdot s^{-3}
\]

\vspace{0.2cm}

On reconnaît immédiatement la signature du joule,
avec \textbf{un facteur $s^{-1}$ en plus}.
Cela confirme que la puissance est bien un \textbf{débit d’énergie} :
des joules par seconde.

\vspace{0.3cm}

\textbf{Lien avec les autres expressions de la puissance}

\vspace{0.1cm}

\[
P = F v
\qquad\text{et}\qquad
P = U I
\]

Ces expressions donnent la puissance sous d’autres formes,
mais l’unité finale reste toujours le watt.

\vspace{0.3cm}

\textbf{Idée clé}

\vspace{0.1cm}

L’énergie est une \textbf{quantité}.  
La puissance est un \textbf{débit}.

Même énergie, puissance différente $\Rightarrow$ durée différente.

}



% -------------------------------------------------
\theorybox{I - Analyse dimensionnelle (signature des unités)}{
Chaque unité possède une \textbf{signature unique}, comme une factorisation en nombres premiers.

\[
J = kg \cdot m^2 \cdot s^{-2}
\]

Si les unités ne correspondent pas, le résultat \textbf{n’est pas une énergie}.  

Exemple :
\[
E = m[kg]\cdot g[ms^{-2}]\cdot h[m]
\Rightarrow
kg \times \frac{m}{s^2} \times m = kg \cdot m^2 \cdot s^{-2}
\]

L’analyse dimensionnelle est un \textbf{outil de détection d’erreur}.
}
% -------------------------------------------------
\theorybox{J - Formules courantes d'énergie et de puissance}{
\vspace{0.2cm}

\textbf{Énergies courantes}

\begin{center}
\begin{tabular}{|>{\small}l|>{\small}l|>{\small}l|>{\small}l|}
\hline
\textbf{Nom} & \textbf{Formule} & \textbf{Paramètres} & \textbf{Signature SI} \\
\hline
Potentielle (gravité) & $E_p = m g h$ & $m$ masse [kg], $g$ gravité [m/s²], $h$ hauteur [m] & $kg \cdot m^2 \cdot s^{-2}$ \\
\hline
Cinétique & $E_c = \frac{1}{2} m v^2$ & $m$ masse [kg], $v$ vitesse [m/s] & $kg \cdot m^2 \cdot s^{-2}$ \\
\hline
Thermique & $E_{th} = m c \Delta T$ & $m$ masse [kg], $c$ chaleur spé... [J/kg·K], $\Delta T$ [K] & $kg \cdot m^2 \cdot s^{-2}$ \\
\hline
Travail d'une force & $W = F d$ & $F$ force [N], $d$ distance [m] & $kg \cdot m^2 \cdot s^{-2}$ \\
\hline
Électrique & $E = U I t$ & $U$ tension [V], $I$ courant [A], $t$ durée [s] & $kg \cdot m^2 \cdot s^{-2}$ \\
\hline
Changement d'état & $E_\text{ch} = m L$ & $m$ masse [kg], $L$ chaleur latente [J/kg] & $kg \cdot m^2 \cdot s^{-2}$ \\
\hline
\end{tabular}
\end{center}

\vspace{0.3cm}

\textbf{Puissances courantes}

\begin{center}
\begin{tabular}{|l|l|l|l|}
\hline
\textbf{Nom} & \textbf{Formule} & \textbf{Paramètres} & \textbf{Signature SI} \\
\hline
Générale & $P = \frac{E}{t}$ & $E$ énergie [J], $t$ temps [s] & $kg \cdot m^2 \cdot s^{-3}$ \\
\hline
Mécanique (force) & $P = F v$ & $F$ force [N], $v$ vitesse [m/s] & $kg \cdot m^2 \cdot s^{-3}$ \\
\hline
Électrique & $P = U I$ & $U$ tension [V], $I$ courant [A] & $kg \cdot m^2 \cdot s^{-3}$ \\
\hline
\end{tabular}
\end{center}

\vspace{0.3cm}

\textbf{Note :}  
La signature SI de toutes les énergies est toujours $kg \cdot m^2 \cdot s^{-2}$,  
et celle de toutes les puissances est $kg \cdot m^2 \cdot s^{-3}$,  
ce qui correspond à des joules par seconde.
}


% =================================================
% PARTIE II : EXERCICES
% =================================================

\newpage
\section*{II - Exercices}

% -------------------------------------------------
\subsection*{Exercice 1 - Train et énergie potentielle}

Un train de masse $400\,\si{t}$ monte une colline de $80\,\si{m}$.

\begin{itemize}
\item Faire un dessin de la situation avec les données connues.
\item Convertir toutes les données en unités SI.
\item Calculer l’énergie potentielle gagnée.
\end{itemize}

\fullgrid

% -------------------------------------------------
\newpage
\subsection*{Exercice 2 - Énergie cinétique d’une voiture}

Une voiture de $1200\,\si{kg}$ roule à $120\,\si{km/h}$.

\begin{itemize}
\item Faire un dessin de la situation avec les données connues.
\item Convertir la vitesse en m/s.
\item Calculer l’énergie cinétique.
\end{itemize}

\fullgrid

% -------------------------------------------------
\newpage
\subsection*{Exercice 3 - Vélo en descente}

Un vélo de 100 kg descend une colline de $25\,\si{m}$ de hauteur. Calculer la vitesse au bas de la colline en km/h. On négligera les frottements.

\begin{itemize}
\item Faire un dessin de la situation avec les données connues.
\item Calculer l'énergie.
\item Calculer la vitesse au bas de la colline.
\item Convertir la vitesse dans l'unité demandée.
\end{itemize}

\fullgrid

% -------------------------------------------------
\newpage
\subsection*{Exercice 4 - Vélo en monté}

Un vélo arrive au pied d'une colline avec une vitesse de 120 km/h. Quelle est la hauteur de la colline sachant qu'il est ralentie en montant et il s'arrête tout juste au sommet. 

\begin{itemize}
\item Faire un dessin de la situation avec les données connues.
\item Convertir la vitesse en m/s.
\item Calculer l'énergie.
\item Calculer la hauteur.
\end{itemize}

\fullgrid

% -------------------------------------------------
\newpage
\subsection*{Exercice 5 - Ascenseur}

Un ascenseur de masse totale $900\,\si{kg}$ est tiré par un moteur de $1000\,\si{W}$ et doit monter une tour de $120\,\si{m}.

\begin{itemize}
\item Faire un dessin de la situation avec les données connues.
\item Calculer l’énergie total.
\item Calculer le temps de montée avec la puissance.
\end{itemize}

\fullgrid

% -------------------------------------------------
\newpage
\subsection*{Exercice 6 - Barrage hydraulique}

Un barrage laisse passer $2000\,\si{kg/s}$ d’eau depuis $600\,\si{m}$. Le lac mesure $2\,\si{km} \times 5\,\si{km}$ et $50\,\si{m}$ de profondeur. On néglige la profondeur du lac pour le calcul. N'utilise que la hauteur de la conduite forcée de $600\,\si{m}$ Calculer la puissance maximale que le barrage peut fournir en Kilo Watt ainsi que l'énergie totale stockée dans le lac en Kilo Watt heure.

\begin{itemize}
\item Faire un dessin de la situation avec les données connues.
\item Convertir en unité SI.
\item Calculer la puissance.
\item Calculer l’énergie stockée dans le lac
\end{itemize}

\fullgrid

% -------------------------------------------------
\newpage
\subsection*{Exercice 7 - Batterie}

Une batterie fournit $24\,\si{V}$ et $15\,\si{A}$ pendant $3\,\si{h}$. Calculer la puissance puis calculer l'énergie total en Joule ainsi qu'en Watt Heure.

\begin{itemize}
\item Faire un dessin de la situation avec les données connues.
\item Convertir en unité SI.
\item Calculer la puissance.
\item Calculer l’énergie en joules
\item Convertir en Wh.
\end{itemize}

\fullgrid

% -------------------------------------------------
\newpage
\subsection*{Exercice 8 - Chauffage}

Un radiateur de $2000\,\si{W}$ fonctionne pendant $45\,\si{min}$.

\begin{itemize}
\item Faire un dessin de la situation avec les données connues.
\item Convertir en unité SI.
\item Calculer l’énergie consommée.
\end{itemize}

\fullgrid

% -------------------------------------------------
\newpage
\subsection*{Exercice 9 - Saut vertical}

Une personne de $70\,\si{kg}$ saute verticalement de $40\,\si{cm}$. Quelle est sa vitesse en retombant juste avant de toucher le sol ?

\begin{itemize}
\item Faire un dessin de la situation avec les données connues.
\item Convertir en unité SI.
\item Calculer l’énergie potentielle.
\item En déduire la vitesse.
\end{itemize}

\fullgrid

% -------------------------------------------------
\newpage
\subsection*{Exercice 10 - Turbine}

Une turbine reçoit $30'000\,\si{litres/s}$ d’eau depuis $25\,\si{m}$. Quelle est la puissance théorique disponible sur l'arbre de la turbine ?

\begin{itemize}
\item Faire un dessin de la situation avec les données connues.
\item Convertir en unité SI.
\item Calculer la puissance.
\end{itemize}

\fullgrid

% -------------------------------------------------

\newpage
\subsection*{Exercice 12 - Panneau solaire}

Des panneaux solaires de surface $12\,\si{m^2}$ sont exposés à un soleil intense
(par beau temps, soleil haut dans le ciel).

L’irradiance solaire au sol est $ I = 1000\,\si{W/m^2} $

Le panneau a une efficacité de $25\%$.
Calcul la puissance électrique disponible à la sortie du panneau en [W] ainsi que la quantité d'énergie électrique total produite entre 11h et 13h en [kWh].
Rappel : Une efficacité de $25\%$ signifie que \textbf{seulement un quart} de l’énergie lumineuse est transformée en électricité. Le reste est principalement perdu sous forme de chaleur.

\begin{itemize}
\item Faire un schéma simple du panneau et du rayonnement solaire.
\item Identifier les données connues et leurs unités.
\item Calculer la puissance solaire reçue par le panneau.
\item Calculer la puissance électrique disponible en tenant compte de l’efficacité.
\item Calculer l’énergie produite.
\end{itemize}


\threequartergrid

% -------------------------------------------------
\newpage
\subsection*{Exercice 13 - Bouilloires et puissance}

Deux bouilloires électriques sont utilisées pour chauffer de l’eau :

\begin{itemize}
\item Bouilloire A : $100\,\si{ml}$ d’eau à $50\,\si{°C}$ avec une puissance de $100\,\si{W}$.
\item Bouilloire B : $1.5\,\si{L}$ d’eau à $20\,\si{°C}$ avec une puissance de $2000\,\si{W}$.
\end{itemize}
la chaleur spécifique de l'eau est $c = 4180\,\si{J/(kg \cdot K)}$.

\textbf{Question :} Laquelle des deux bouilloires portera son eau à $100\,\si{°C}$ en premier et combien de temps faudra-t-il pour chacune ?

\begin{itemize}
\item Faire un dessin de la situation avec les données connues.
\item Convertir toutes les données en unités SI :
\item Calculer l’énergie nécessaire pour chaque cas.
\item Calculer le temps à partir de la puissance
\end{itemize}

\threequartergrid

\end{document}
